\subsection*{Chapter 2: Equilibrium Constant of Reaction}

\subsubsection*{Redox Reactions}
\begin{itemize}
    \item \textbf{Big positive E} = \textbf{stronger oxidizing} agent and is more \textbf{easily reduced} = \textbf{reduction half reaction}.
    \item \textbf{n}: number of mol e transferred per mole of that substance according to its \textbf{half-reaction}.
    \item \begin{itemize}
    \item \textbf{Redox (Oxidation-Reduction):}
    \begin{itemize}
        \item $n$ is the number of \textbf{electrons ($\mathrm{e}^-$)} gained or lost.
        \item \textit{Example:} For $\mathrm{MnO}_4^{-} \rightarrow \mathrm{Mn}^{2+}$, $n=5$.
    \end{itemize}

    \item 
    \textbf{Acid-Base (Neutralization):}
    \begin{itemize}
        \item $n$ is the number of \textbf{transferable protons ($\mathrm{H}^{+}$)} or $\mathrm{OH}^{-}$ ions/acceptor sites.
        \item \textit{Example:} For $\mathrm{H}_3\mathrm{PO}_4$ (fully reacted), $n=3$.
    \end{itemize}

    \item \textbf{Precipitation:}
    \begin{itemize}
        \item $n$ is the \textbf{total magnitude of the charge} contributed by the reacting ion(s) in the formula unit.
        \item \textit{Example:} For $\mathrm{Al}_2(\mathrm{SO}_4)_3$, $n = 6$ (from $2 \times |+3|$ charge).
    \end{itemize}

    \item \textbf{Complexation:}
    \begin{itemize}
        \item $n$ is typically the \textbf{charge of the central metal ion}.
        \item \textit{Example:} For $\mathrm{Ag}^{+}$, $n=1$.
    \end{itemize}
\end{itemize}

\end{itemize}
The value of $\mathbf{n}$ in the equivalent weight expression $\left(E = \frac{M}{n}\right)$ is the number of moles of the fundamental reactive unit transferred per mole of the substance:


\textbf{Nernst Equation for Reduction Potential:}
\begin{equation}
E = E^\circ + \frac{RT}{nF} \ln \frac{[Ox]^a}{[Red]^b}
\end{equation}
\textit{E}: Electrode potential (V), \textit{E°}: Standard electrode potential (V), \textit{R}: Gas constant (8.314 J/mol·K), \textit{T}: Temperature (K), \textit{n}: Number of electrons, \textit{F}: Faraday constant (96493 C/mol), \textit{[Ox], [Red]}: Molarities.

\vspace{5mm}

\textbf{Nernst Equation (simplified at 25°C):}
\begin{equation}
E = E^\circ + \frac{0.059}{n} \log \frac{[Ox]}{[Red]}
\end{equation}

\textbf{Equilibrium Constant for Redox Reaction:}
\begin{equation}
K = 10^{\frac{n_1 n_2 (E_1^\circ - E_2^\circ)}{0.059}}
\end{equation}
\textit{n\textsubscript{1}, n\textsubscript{2}}: Number of electrons for each half-reaction, \textit{E°\textsubscript{1}, E°\textsubscript{2}}: Standard potentials for each half-reaction.

\vspace{5mm}

\textbf{Equivalence Point Potential:}
\begin{equation}
E_{\text{eq}} = \frac{n_1 E_1^\circ + n_2 E_2^\circ}{n_1 + n_2}
\end{equation}

\vspace{5mm}

\textbf{Presence of H\textsuperscript{+} in half-reaction of Ox\textsubscript{1}/Red\textsubscript{1}:}
\begin{equation}
E_{\text{eq}} = \frac{n_1 E_1^\circ + n_2 E_2^\circ}{n_1 + n_2} + \frac{0.059}{n_1 + n_2} \lg[H^+]^m
\end{equation}

\vspace{5mm}

\textbf{Presence of H\textsuperscript{+} and different stoichiometric coefficients of Ox\textsubscript{1} and Red\textsubscript{1}:}

\textbf{1. Reduction Half-Reaction ($\text{Ox}_{1} \rightarrow \text{Red}_{1}$)}
\begin{equation}
\text{Ox}_{1} + n_{1}\text{e}^{-} + m\text{H}^{+} \rightleftharpoons p\text{Red}_{1} + \frac{1}{2}m\text{H}_{2}\text{O}
\end{equation}

\textbf{2. Oxidation Half-Reaction ($\text{Red}_{2} \rightarrow \text{Ox}_{2}$)}
\begin{equation}
\text{Red}_{2} \rightleftharpoons \text{Ox}_{2} + n_{2}\text{e}^{-}
\end{equation}

\textbf{3. Balanced Overall Reaction}


\begin{equation}
n_{2}\text{Ox}_{1} + n_{1}\text{Red}_{2} + n_{2}m\text{H}^{+} \mathrel{\mathop{\rightleftharpoons}\limits_{\text{(2)}}^{\text{(1)}}} n_{1}\text{Ox}_{2} + n_{2}p\text{Red}_{1} + \frac{1}{2}n_{2}m\text{H}_{2}\text{O}
\end{equation}

\textbf{4. Equivalence Point Potential Formula}

\begin{equation}
E_{\text{eq}} = \frac{n_1 E_1^\circ + n_2 E_2^\circ}{n_1 + n_2} + \frac{0.059}{n_1 + n_2} \lg \left[ \frac{[H^+]^m [\text{Red}_1]^{1-p}}{p} \right]
\end{equation}

\subsubsection*{Dissociation and Formation Constants}

\textbf{Stepwise Constants} ($K(1)$, $\beta$, $\beta_i$, $k'_i$)\\
For the reaction \( A + p \rightleftharpoons D \):

\begin{equation}
K(1) = \beta = \frac{[D]}{[A][p]} \quad (\beta (K_f): \text{formation constant})
\end{equation}
For the general "i" step, \( D_{i-1} + p \rightleftharpoons D_i \):
\begin{equation}
\beta_i = \frac{[D_i]}{[D_{i-1}][p]} = \frac{1}{k'_i}
\end{equation}
\textbf{Dissociation Constant:} ($k$)
\begin{equation}
k = \frac{[A][p]}{[D]} = \frac{1}{\beta}
\end{equation}

\vspace{5mm}

\textbf{Overall Formation Constant} ($\beta_{1,i}$)\\
For an n-step equilibrium:
\begin{equation}
\beta_{1,i} = \beta_1 \cdot \beta_2 \cdot \ldots \cdot \beta_i = \frac{1}{k_n \cdot k_{n-1} \cdot \ldots \cdot k'_i}
\end{equation}

\vspace{5mm}

\textbf{Relationship between Overall and Stepwise Constants}\\
For a 2-step equilibrium:
\begin{equation}
\beta_1 \cdot \beta_2 = \frac{[D_1]}{[A][p]} \times \frac{[D_2]}{[D_1][p]} = \frac{[D_2]}{[A][p]^2} = \beta_{1,2}
\end{equation}

\subsubsection*{Concentration at Equilibrium}

\textbf{Mass Balance} ($[A]_o$)
\begin{equation}
[A]_o = [A] + [D_1] + [D_2] + \ldots + [D_n]
\end{equation}

\textbf{Conditional Coefficient} ($\alpha_{A[p]}$ - only once)
\begin{equation}
[A]_o = [A] \cdot \left\{1 + \sum_{i=1}^{n} \beta_{1,i}[p]^i\right\} = [A] \cdot \alpha_{A[p]}
\end{equation}

\textbf{Equilibrium Concentrations} ($[A]$, $[D_i]$)
\begin{equation}
[A] = \frac{[A]_o}{\alpha_{A[p]}}
\end{equation}
\begin{equation}
[D_i] = \frac{[A]_o \beta_{1,i}[p]^i}{\alpha_{A[p]}}
\end{equation}

\vspace{5mm}

\textbf{For EDTA (H\textsubscript{4}Y)}
\begin{equation}
\alpha_{Y(H^+)} = 1 + \beta_{1,1}[H^+] + \beta_{1,2}[H^+]^2 + \beta_{1,3}[H^+]^3 + \beta_{1,4}[H^+]^4
\end{equation}

\textbf{Concentration of EDTA Species}
\begin{align}
[Y^{4-}] &= \frac{[Y]_o}{\alpha_{Y(H^+)}} \\
[HY^{3-}] &= \frac{[Y]_o 10^{10.95} [H^+]}{\alpha_{Y(H^+)}} \\
[H_2Y^{2-}] &= \frac{[Y]_o 10^{17.22} [H^+]^2}{\alpha_{Y(H^+)}} \\
[H_3Y^-] &= \frac{[Y]_o 10^{19.89} [H^+]^3}{\alpha_{Y(H^+)}} \\
[H_4Y] &= \frac{[Y]_o 10^{21.89} [H^+]^4}{\alpha_{Y(H^+)}}
\end{align}

\subsubsection*{Precipitation Reactions} 

\begin{equation}
A + n p \stackrel{\beta_D}{\rightleftharpoons} D \stackrel{\beta_{D\downarrow}}{\rightleftharpoons} D\downarrow
\end{equation}

\textbf{Formation Constant}
\begin{equation}
\beta_D = \frac{[D]}{[A][p]^n}
\end{equation}
\begin{equation}
\beta_{D\downarrow} = \frac{1}{[D]}
\end{equation}

\textbf{Solubility Product (K\textsubscript{sp})}
\begin{equation}
K_{sp} = \frac{1}{\beta_D \cdot \beta_{D\downarrow}}
\end{equation}
For a slightly soluble compound \( A_mB_n \rightleftharpoons mA^{n+} + nB^{m-} \):
\begin{equation}
K_{sp} = [A^{n+}]^m[B^{m-}]^n
\end{equation}
If the salt $MA$ has the anion $A^-$ which is the conjugate base of the weak acid $HA$:
$${MA + H_2O \rightleftharpoons M+ + HA + OH-} $$
then:
$$ K_{\text{total}} = K_{sp} \cdot K_a $$

and if the pH is known:

$$ S = \sqrt{\frac{K_{\text{total}}}{[OH^-]}} = \sqrt{\frac{K_{sp} \cdot K_a}{[OH^-]}} $$

\textbf{Solubility (S)}
\begin{equation}
S = [D] + [A]
\end{equation}
\begin{equation}
S = \sqrt[m+n]{\frac{Ksp_{A_mB_n}}{m^m \cdot n^n}}
\end{equation}
For Ag\textsubscript{2}CrO\textsubscript{4}:
\begin{equation}
K_{sp} = (2s)^2 \cdot s = 4s^3
\end{equation}

\subsubsection*{Acid-Base Reactions}
\begin{equation}
    [H^+] = 10^{-pH}
\end{equation}

\textbf{Ionic Product of Water (K\textsubscript{w})}
\begin{equation}
k_w = [H^+][OH^-]
\end{equation}
At 25°C, \( k_w = 10^{-14} \).
\begin{equation}
pH + pOH = pk_w = 14
\end{equation}

\textbf{Acid and Base Dissociation Constants}
\begin{equation}
k_a = \frac{[H^+][A^-]}{[HA]}
\end{equation}
\begin{equation}
k_b = \frac{[HA][OH^-]}{[A^-]} = \frac{k_w}{k_a}
\end{equation}

\subsubsection*{Acid-Base pH Calculations}

\textbf{pH of Monoprotic Acids}
\begin{itemize}
    \item \textbf{Strong Acid ($k > 10^0$):}
    \begin{itemize}
        \item When \( [HA] > 10^{-7} \) M:
        \begin{equation}
        pH = -\log([HA])
        \end{equation}
        \item When \([HA] \sim 10^{-7} \) M (dissociation of water is not negligible):
        \begin{equation}
        [H^+]^2 - [HA][H^+] - 10^{-14} = 0
        \end{equation}
    \end{itemize}
    \item \textbf{Medium Acid ($10^{-3} < k < 10^0$):}
    \begin{equation}
    [H^+]^2 + k_a[H^+] - k_a[HA] = 0
    \end{equation}
    \item \textbf{Weak Acid ($k < 10^{-3}$):}
    \begin{equation}
    \text{p}k_a = -\log(k_a)
    \end{equation}
    \begin{equation}
    pH = \frac{1}{2} pk_a - \frac{1}{2} \log([HA])
    \end{equation}
\end{itemize}

\textbf{Case: If another monoprotic acid is added:}
\begin{equation}
n_a = M_aV_a
\end{equation}
\begin{equation}
pH = pk_a + \log \frac{[A^-] - n_{a}}{[HA] + n_{a}}
\end{equation}
\textbf{$n_{a}$}: moles of added monoprotic acid.
\textbf{$M_{a}$}: Molar conc. of added monoprotic acid.
\textbf{$V_{a}$}: Volume of added monoprotic acid. (L)

\textit{Note: If [HA] isn't specified, [HA] = 0}
\vspace{5mm}

\textbf{pH of Polyprotic Acid / Mixture of Monoprotic Acids}
\begin{itemize}
    \item If dissociation constants are \textbf{very different} (\( k_1 \gg k_2 \), e.g., by a factor of \( \sim 10^4 \)):
    \begin{itemize}
        \item The pH is determined primarily by the first dissociation. Treat as a monoprotic acid with constant \( k_1 \).
    \end{itemize}
    \item If dissociation constants are \textbf{similar} (\( k_1 \sim k_2 \)):
    \begin{equation}
    [H^+]^2 = \sum_{i=1}^{n} k_{HA_i} \cdot [HA_i]
    \end{equation}
\end{itemize}

\vspace{5mm}

\textbf{pH of Monoprotic Bases}
\begin{itemize}
    \item \textbf{Strong Base:}
    \begin{itemize}
        \item When \( [B] > 10^{-7} \) M:
        \begin{equation}
        pH = 14 + \log[OH^-]
        \end{equation}
        \item When \( [B] \sim 10^{-7} \) M (dissociation of water is not negligible):
        \begin{equation}
        [OH^-]^2 - [B][OH^-] - 10^{-14} = 0
        \end{equation}
    \end{itemize}
    \item \textbf{Medium Base:}
    \begin{equation}
    [OH^-]^2 + k_b[OH^-] - k_b[B] = 0
    \end{equation}
    \item \textbf{Weak Base:}
    \begin{equation}
    \text{p}k_b = -\log(k_b)
    \end{equation}
    \begin{equation}
    pH = 14 - \frac{1}{2} pk_b + \frac{1}{2} \log([B])
    \end{equation}
\end{itemize}

\vspace{5mm}

\textbf{pH of Polyprotic Base / Mixture of Monoprotic Bases}
\begin{itemize}
    \item If dissociation constants are \textbf{very different} (\( k_1 \gg k_2 \), e.g., by a factor of \( \sim 10^4 \)):
     \begin{itemize}
        \item The pH is determined primarily by the first dissociation. Treat as a monoprotic base with constant \( k_1 \).
    \end{itemize}
    \item If dissociation constants are \textbf{similar} (\( k_1 \sim k_2 \)):
    \begin{equation}
    [OH^-]^2 = \sum_{i=1}^{n} k_{B_i} \cdot C_{B_i}
    \end{equation}
\end{itemize}

\vspace{5mm}

\textbf{pH of Buffer Solutions (Henderson-Hasselbalch Equation)}
\begin{itemize}
    \item A buffer is a mixture of a weak acid and its conjugate base, or a weak base and its conjugate acid.
    \begin{equation}
    pH = pk_a + \log\left(\frac{[\text{Base}]}{[\text{Acid}]}\right) \quad \text{or} \quad [H^+] = k_a \cdot \frac{C_A}{[B]}
    \end{equation}
\end{itemize}

\vspace{5mm}

\textbf{pH of Salt Solutions}
\begin{itemize}
    \item \textbf{Salt from Strong Acid and Strong Base} (e.g., NaCl):
    \begin{itemize}
        \item The solution is neutral, \( pH = 7 \) at 25°C.
    \end{itemize}
    \item \textbf{Salt from Strong Acid and Weak Base} (e.g., NH\textsubscript{4}Cl):
     \begin{itemize}
        \item The solution is acidic.
        \begin{equation}
        pH = \frac{1}{2} pk_a - \frac{1}{2} \log(C_{salt})
        \end{equation}
    \end{itemize}
    \item \textbf{Salt from Weak Acid and Strong Base} (e.g., CH\textsubscript{3}COONa):
     \begin{itemize}
        \item The solution is basic.
        \begin{equation}
        pH = 7 + \frac{1}{2} pk_a + \frac{1}{2} \log(C_{salt})
        \end{equation}
    \end{itemize}
\end{itemize}

